\section{Conclusion and Future Work}
\label{sec:future_work}

The analysis of the consistency- and coherence-lattices highlights the distinction between strict and relaxed notions of agreement between knowledge bases. 
The consistency-lattice retains only concepts whose extents coincide exactly in both the \ac{mlb}-constrained and the topic model-based lattices. 
In our experiments, no non-trivial document clusters were shared exactly between the two lattices. %, indicating that strict consistency is often too restrictive when comparing knowledge bases derived from heterogeneous sources.
% limitation
Generally, pruning low-support concepts through iceberg filtering can remove exact overlaps and thus yield an incomplete consistency-lattice. 
Thus, it is important to design the pipeline data-dependent to avoid this issue.
% However, in our experiments, the absence of non-trivial shared concepts appears to reflect genuine semantic differences between the expert-defined and topic model-based knowledge bases rather than an artefact of pruning.

The empty consistency-lattice refelcts fundamental differences between the two knowledge bases.
The \ac{mlb}-constrained lattice is derived from manual annotations and thus reflects human expertise, preferences and biases, while the topic model-based lattice is data-driven and captures latent statistical regularities and noise in the corpus, potentially revealing structures not explicitly encoded in the manual taxonomy.
Our approach allows for the integration of alternative topic modeling methods.

Our study provides two main contributions. 
First, we showed that \ac{mlb} constraints can be integrated natively into the \ac{fca} framework through a fixpoint-based closure operator. 
Second, experiments on the BankSearch dataset demonstrated that exact consistency between manual hierarchies and data-driven topic models is rarely observed in real-world corpora.

In practice, this implies that a multi-tier approach is necessary: (1) using the \textit{consistency-lattice} to identify undisputed common ground and (2) utilizing the \textit{coherence-lattice} with tunable thresholds to explore meaningful alignments. 

The coherence-lattice addresses this limitation by building document clusters based on partial overlaps between concept extents through a similarity threshold.  
Moreover, the number of documents contained in coherence-lattice concepts provides an indication of the degree of agreement between the underlying representations.

Overall, the results suggest that practical knowledge integration requires tolerating a certain degree of inconsistency while maintaining sufficient coherence. 
The appropriate similarity threshold depends on the application domain: safety-critical settings such as medicine may require high thresholds to ensure reliable alignment, whereas exploratory analysis may benefit from lower thresholds that permit broader, though less precise, overlaps.

Future work will investigate the integration of more than two knowledge bases, methods for selecting similarity thresholds that yield coherence-lattices of comparable size, and alternative overlap measures beyond Jaccard similarity for comparing concept extents across multiple lattices.